\section{Introduction}\label{sec:intro}

Equational reasoning is the daily bread of circuit design.  Whether one
optimizes a quantum circuit, verifies a compiler pass, or normalizes a
diagram, the justification is always the same: a finite set of local rewrite
rules, applied to a syntactic object, preserving a semantic interpretation.
Two questions attach to any such rule set.  \emph{Soundness} --- does every
rule preserve the semantics? --- is typically easy: it is a finite collection
of small algebraic identities.  \emph{Completeness} --- can every true
semantic equation be derived from the rules? --- is a different animal
altogether.

Completeness proofs are gigantic computations.  The standard strategy defines
a \emph{normal form} for circuits, proves that every circuit can be rewritten
to normal form using the rules, and proves that distinct normal forms denote
distinct semantic objects.  The middle step is where the explosion happens:
one must show, for every rule and every way a generator can meet a normal
form, that the combination reduces again --- a case analysis that is
mechanical, unforgiving, and enormous.  The literature wears the scars
openly.  Selinger's presentation of the $n$-qubit Clifford group already
required a substantial rewrite system and normal form
analysis~\cite{selinger2015clifford}.  Its recent qutrit
analogue by Li, Mosca, Ross, van~de~Wetering, and Zhao --- the first
completeness result for a fragment of quantum circuits in an odd prime
dimension --- runs to more than fifty pages of
proofs~\cite{li2025qutrit}.  The first complete equational theory for full
quantum circuits, by Cl\'ement, Heurtel, Mansfield, Perdrix, and
Valiron~\cite{clement2023complete}, and its minimal
refinement~\cite{clement2024minimal}, rest on intricate normal-form
manipulations; and Cl\'ement's complete equational theory for
Real-Clifford+$CH$ circuits consists of a thirteen-page paper attached to an
appendix of \emph{264 pages} of equational
derivations~\cite{clement2026realcliffordch}.  These are not outliers: the
same tens of pages of mechanical calculation stand behind Selinger's
normal-form analysis and its qutrit successor cited above, behind the
rig-categorical completeness theorems for quantum programming
languages~\cite{carette2024squareroots,fang2026hadamard}, and behind the
recent uniform simplification of the rule sets of six near-Clifford circuit
fragments~\cite{blake2026simpler} --- long, hand-written, hand-checked
derivations every one of which deserves a machine's scrutiny.  Proofs of
this shape --- long, regular, and computational --- are precisely the proofs
that should be machine-checked, and precisely the proofs that a proof assistant with a
strong evaluator can check largely \emph{by computation}.

This paper reports on an Agda library built for exactly this purpose, and
on the design decisions that make it work.  The library is a general,
reusable framework for \emph{presentations of monoids and groups by
generators and relations}, for \emph{normal forms}, and for
\emph{completeness} of finite relation sets, with circuits --- words of
wire-indexed gates --- as the motivating instance.  Its approach can be
stated without formal vocabulary.  A circuit is a word of gates, and a rule
set is a datatype of axioms; the equational theory is the congruence these
axioms generate, itself an inductive datatype, so that a derivation is a
first-class object.  A normal form is an ordinary function on words, paired
with a function back from normal forms to words that picks a canonical
representative; completeness then follows from soundness together with one
further property, that normal forms with equal denotations are equal ---
and this reduction is proved once, as a generic lemma, never per example.
Normal forms themselves are built one wire at a time.  The circuits on $n$
wires sit inside those on $n{+}1$ wires as the circuits that leave the
bottom wire alone, and a \emph{coset table} --- a small first-order function
recording how a gate moves past the representative of each coset of that
subgroup --- adds one ``digit'' to the normal form per wire.  The table is
not searched for inside the proof assistant; it is supplied, and
\emph{verified} by checking five equations.  Because tables and normal
forms are structurally recursive programs, those equations, and the
hundreds of case-split obligations they expand into, are closed by Agda's
evaluator: the case explosion still happens, but inside the typechecker.
Section~\ref{sec:permutations} shows the whole route on the smallest
example, permutation circuits presenting the symmetric groups.

Two design decisions deserve to be announced up front, since the paper
returns to them repeatedly.  First, the library is written in plain,
setoid-based Agda on top of the standard
library~\cite{daggitt2025stdlib,agdastdlib}, not in cubical Agda or a
HoTT library: a presented monoid is a setoid on words, not a quotient type.
The reason is that our central devices are maps \emph{into} the syntax ---
sections of normal forms, coset representatives, amalgamation tails ---
which are natural for setoids and awkward for quotients;
\S\ref{sec:presentations} and \S\ref{sec:setoids} weigh the trade.
Second, the library verifies rather than searches: coset tables are found
offline and certified inside Agda, which keeps the obligations first-order
and the trusted story simple (\S\ref{sec:lessons}).

\paragraph{Why this is a programming-languages problem.}
Circuit calculi are programming languages of a particular kind, and
completeness of their equational theories is a programming-languages
question: it is what licenses normal-form-based decision procedures, what
tells the author of an optimizer that no valid rewrite is missing, and what
underlies the rig-categorical completeness theorems for reversible and
quantum languages of recent
POPLs~\cite{choudhury2022symmetries,carette2024squareroots,fang2026hadamard}.
The contribution of this paper is not a single theorem but a \emph{design}:
a way of organizing presentations, normal forms, and completeness in a
dependently typed language so that the same objects both compute and
prove, validated on examples large enough --- the qutrit Clifford+$T$
monoid, $U_3(\mathbb{Z}[\tfrac12,i])$ --- to exercise every part of it.
The techniques that carry the design are this community's: inductive
congruences, specifications as function bundles, index engineering for
the unifier, and proof by evaluation.

\subsection{Contributions}\label{sec:contributions}

\begin{enumerate}
\item[(0)] \textbf{A library for circuit normal forms and complete relation
  sets.}  An Agda library, organized in layers, in which a relation set is
  an inductive family on words and the presented monoid is the inductively
  generated congruence closure (\S\ref{sec:presentations}); normal-form
  witnesses are the standard library's function bundles out of the setoid
  of words (\S\ref{sec:nfproperty}); and a single generic lemma,
  \aid{by-normalization}, turns soundness plus a \emph{unique normal form}
  into completeness.  A verified, monoid-level Reidemeister--Schreier
  engine constructs normal forms from coset tables, one verified table per
  level of a subgroup tower (\S\ref{sec:engine}).

\item[(1)] \textbf{Inductively defined circuits, for easy normalization and
  completeness proofs.}  A generic circuit layer represents $n$-wire circuits
  as words over wire-indexed generators, and extends any gate-specific
  relation set with the structural rules --- shift congruence and
  far-commutation --- shared by every circuit presentation
  (\S\ref{sec:circuits}).  The indexing discipline is chosen so that Agda's
  evaluator, rather than the human, carries out the case analyses: in our
  case studies, every coset-table obligation closes by \aid{refl} once both
  sides have computed.

\item[(2)] \textbf{A catalogue of verified presentations}
  (\S\ref{sec:examples}).  Seven nontrivial example families, several of them
  formalizing presentations from the research literature: cyclic groups
  of every positive order (with, at order $0$, the free monoid
  $\mathbb{N}$ presented as a monoid --- the smallest witness that our
  presentations are monoid presentations); symmetric groups $S_n$ presented as circuits,
  proved sound and complete for a loose endofunction semantics and presented
  for the group of permutations of $\mathsf{Fin}\;n$; wreath products
  $\mathbb{Z}_N \wr S_n$ --- the groups of \emph{signed permutations} and
  their generalizations; Pauli groups $(\mathbb{Z}_p \times \mathbb{Z}_p)^n$
  for every odd prime $p$, assembled purely compositionally; and, as
  two- and four-level amalgamated free products with verified alternating
  normal forms, the single-qubit Clifford+$T$ monoid, a single-qutrit
  Clifford+$T$ monoid, and $U_3(\mathbb{Z}[\tfrac12,i])$ from Bian and
  Selinger's presentation of
  $U_n(\mathbb{Z}[\tfrac12,i])$~\cite{bian2021un}.  Every headline theorem is
  restated, with its full type, in a single index module
  \texttt{MainTheorems.agda}.

\item[(3)] \textbf{Monoid-level amalgamation and Reidemeister--Schreier.}
  Following the methodology that Bian and Selinger introduced for their
  syntactic verifications~\cite{bian2023cliffordt2,bian2023cliffordcs3,
  bian2023thesis}, we generalize the group-theoretic notions ---
  coset enumeration, the Reidemeister--Schreier subgroup presentation, and
  amalgamated free products --- from groups to monoid presentations, where
  inverses are not primitive.  Beyond the syntactic side, we prove the
  corresponding \emph{semantic} presentation theorems: if $\Gamma$ presents
  $G_1$ and $\Delta$ presents $G_2$, the constructed relation sets present
  the direct product, the semidirect product, the $n$-fold product, the
  amalgamated free product, and (given a realizing short exact sequence) an
  arbitrary group extension (\S\ref{sec:constructions}).

\item[(4)] \textbf{New theorems for the standard library.}  The semantic
  side of the above required algebraic constructions absent from the Agda
  standard library: semidirect products of monoids and groups via a bundled
  \aid{Action} record; amalgamated free products of monoids as a pushout
  presented by words over the disjoint union of carriers; group extensions
  and split extensions as short exact sequences; monoid epimorphisms; the
  fact that a monoid homomorphism between groups is a group homomorphism;
  and the group of permutations of $\mathsf{Fin}\;n$ bundled as a
  \aid{Group} (\S\ref{sec:forstdlib}).
\end{enumerate}

Everything stated as a theorem in this paper typechecks under
\texttt{--safe} and \texttt{--cubical-compatible} (the discipline the
standard library calls \texttt{--without-K} again from its version 3.0):
there are no postulates, no unfinished holes, and no appeals to classical
axioms anywhere in the development described here.  Wherever a statement
in this paper is \emph{not} formalized, the text says so explicitly, and
\S\ref{sec:threats} collects these places.  The driving next target ---
machine-checked completeness for multi-qudit Clifford circuits in all odd
prime dimensions, the formal counterpart of the qutrit rule set
of~\cite{li2025qutrit} and its odd-prime generalization --- is discussed
as future work (\S\ref{sec:conclusion}).

\subsection{Related work in brief}\label{sec:related-intro}

The closest prior work is the line of Agda verifications accompanying Bian
and Selinger's presentations of quantum gate
groups~\cite{bian2023cliffordt2,bian2023cliffordcs3,bian2023thesis}.  Those
developments verify that two \emph{syntactic} presentations of a group are
equivalent --- an isomorphism of finitely presented monoids obtained by
Reidemeister--Schreier steps --- and never mention a semantics, so
soundness and completeness with respect to an interpretation are not
stated.  The present library grew out of that methodology and subsumes it:
our amalgamation case studies end in exactly such syntactic isomorphisms,
but the framework additionally relates presentations to independently
constructed semantic groups, which is what licenses the words ``sound''
and ``complete''.  To our knowledge the only other mechanized completeness
results in the vicinity are categorical: the univalent presentation of
the reversible language $\Pi$ as a free rig
groupoid~\cite{choudhury2022symmetries} and the partial formalization ---
complete for the two-qubit Clifford fragment --- accompanying its quantum
square-root extension~\cite{carette2024squareroots}.  The quantum-circuit
verification ecosystems in Rocq and Isabelle verify programs, compilers,
or the soundness of rewrites, not completeness; and the group-theory
libraries of Lean, Rocq, and Isabelle develop presentation theory in the
abstract, without infrastructure for proving that a concrete relation set
presents a concrete group.  Section~\ref{sec:related} makes all of these
comparisons in detail.

\subsection{Paper outline}

Section~\ref{sec:permutations} walks through the permutation case study.
Section~\ref{sec:preliminaries} recalls the group theory behind the design:
normal forms from subgroup chains and transversals, and their mixed-radix
reading.  Section~\ref{sec:design} presents the library layer by layer,
with the design rationale first: words and presentations, the zoo of
normal-form witnesses, inductively defined circuits, the
Reidemeister--Schreier engine, and the product constructions with their
presentation theorems.  Section~\ref{sec:examples} is the catalogue of
verified results and what each was chosen to exercise, with pointers into
\texttt{MainTheorems.agda}.  Section~\ref{sec:related} compares with other
formalizations.  Section~\ref{sec:discussion} discusses statistics and
typechecking cost, design decisions, the choice of Agda, lessons learned,
and limitations, and Section~\ref{sec:conclusion} concludes.
